\section{Cyprian's School of Mortal Life}

St.~Cyprian addresses a community experiencing widespread death in his treatise \emph{On Mortality}. Disease has reached believers as well as their neighbors. Some are unsettled that faith has not protected them from the common affliction. Cyprian answers that Christians share the bodily condition of the human race while awaiting incorruption (no.~8). Believers and unbelievers alike become sick, and the faithful must endure bodily suffering until the promised transformation. He directs their hope toward the life to come while examining how they respond to the epidemic before them.

He recalls the struggle with avarice, lust, anger, ambition, pride, and envy (nos.~4--5). The longing to be with Christ includes the longing to be free of this warfare, to love without the instability that wounds love here. His language is ardent: the Christian is a traveler whose home lies with God. At the close of the treatise he gathers the apostles, prophets, martyrs, virgins, and merciful into the expectation of a heavenly homeland (no.~26). The end of life is contemplated through the company into which the faithful hope to be received.

This desire has an apostolic pattern. Paul longs to depart and be with Christ, yet recognizes that continuing to live serves the needs of the faithful (Phil 1:20--26). Charity holds together the desire for the Lord and the willingness to remain at work. We may ask for healing and more time, especially for the sake of responsibilities entrusted to us. We may also consent to the Lord's summons when it comes. The desire for heaven does not give us authority to hasten death; it orders the heart toward God while leaving our time in his hands.

Cyprian's most searching account of the epidemic concerns what it reveals about the living. Do the healthy tend the sick? Do relatives love their own? Do physicians remain with their patients? Do the powerful show pity, and do the wealthy become generous (no.~16)? Mortality tests more than the composure of the person who dies. It tests the charity of everyone near the suffering. The Christian who remembers death must be ready to spend life in service. An eloquent desire for heaven is hollow if it leaves the sick neighbor without help.

He also addresses believers who fear that an ordinary death has deprived them of the martyrdom for which they had hoped (no.~17). Martyrdom, he answers, is not under their control. God knows the heart's readiness and judges it truly. This passage releases us from trying to choose a spiritually impressive ending. Our task is fidelity to the circumstances God actually permits. A hidden death after long frailty can receive the same wholehearted consent that had imagined itself answering a public trial. The Lord sees the will that illness prevents from displaying its strength.

Cyprian rebukes grief severely, particularly where he fears that Christians mourn as though the departed had been extinguished (nos.~20--21). He urges them to express their hope in the way they mourn. The Gospel also recounts Christ's tears at Lazarus's tomb, and Augustine describes his grief for Monica. These accounts show how deeply the separation can hurt even those who hope in resurrection. We can acknowledge the pain and continue praying for the departed, trusting that death has not destroyed the people whose absence we feel.

Cyprian's questions require us to examine our conduct toward those who are suffering. As we ask God to free us from avarice and ambition, we can make time or possessions available to someone who needs them. A visit, a meal, relief for an exhausted caregiver, or faithful professional care responds to the needs he names. Such service may interrupt the plans we had made for prayer or reflection. It is itself part of Christian preparation, because we are practicing the charity in which we hope to persevere until death.

